% TOP_PROBLEM: {"problemId": "problem.top500-omission-w045049-balanced-trilinear-hilbert-transform", "canonicalTitle": "Balanced L4 boundedness of the trilinear Hilbert transform", "releaseRank": 444, "primaryDomain": "analysis_pde", "primaryDomainLabel": "Analysis & PDE", "displayStatus": "open", "releaseStatus": "open"}
% !TeX program = lualatex
% ATTEMPT_STATUS: unresolved
\documentclass[11pt]{article}
\usepackage[margin=25mm]{geometry}
\usepackage{fontspec}
\setmainfont{Latin Modern Roman}
\usepackage{amsmath,amssymb,mathtools}
\usepackage{unicode-math}
\setmathfont{Latin Modern Math}
\usepackage{xurl,hyperref}
\hypersetup{colorlinks=true,urlcolor=blue}
\setcounter{secnumdepth}{0}
\setlength{\parindent}{0pt}
\setlength{\parskip}{5pt}
\setlength{\emergencystretch}{3em}
\begin{document}
\section*{\texorpdfstring{Balanced L4 boundedness of the trilinear Hilbert transform}{Balanced L4 boundedness of the trilinear Hilbert transform}}\label{444}
\subsection{Definitions and mathematical statement}
For Schwartz functions $f,g,h:{\mathbb{R}}\to{\mathbb{C}}$, define the trilinear Hilbert transform with fixed slopes $1,2,3$ by
\[
 T(f,g,h)(x)=\operatorname{p.v.}\int_{{\mathbb{R}}}
 f(x+t)g(x+2t)h(x+3t)\,\frac{{\,\mathrm{d}} t}{t}.
\]
Principal value means symmetric deletion of an interval around zero. The selected boundedness question is the existence of an absolute finite $C$ with
\[
 \|T(f,g,h)\|_{L^{4/3}({\mathbb{R}})}
 \le C\|f\|_{L^4({\mathbb{R}})}\|g\|_{L^4({\mathbb{R}})}\|h\|_{L^4({\mathbb{R}})}.
\]
The estimate is required uniformly over all such inputs and would give a bounded extension to the indicated spaces. The exponent identity $3/4=1/4+1/4+1/4$ fixes the balanced scaling; this entry does not ask for the entire possible exponent range.
\par\smallskip\textit{Formulation and concept references:} \href{https://terrytao.wordpress.com/2007/05/10/open-question-boundedness-of-the-trilinear-hilbert-transform/}{[S1]}.

\subsection{Short English statement}
Is the trilinear Hilbert transform bounded from three copies of L-four into L-four-thirds at the specified slopes?
\subsection{Research attempt}
\paragraph{Process and selected approach.}
This attempt was generated by GPT-6 Astra Ultra on 27 September 2026. We examine the exact balanced estimate through truncation, Fourier localization, and modulation identities. We prove a uniform estimate on a restricted frequency class, derive the full multiplier with its sign and normalization, and identify a quadratic modulation symmetry that obstructs a naive oscillation argument. None of these calculations gives the requested bound for arbitrary Schwartz inputs.

The first remaining gap is a bound independent of the number of interacting scales. The formulation source [S1] distinguishes this trilinear problem from the bilinear Hilbert transform. Tao's later cancellation theorem [E1] improves the elementary growth with the truncation range but does not establish the full boundedness assertion. Results for curved parameterizations [E2] concern different operators: the three shifts here are linear in the same parameter, with slopes exactly $1,2,3$.

\paragraph{The absolute-value estimate loses the number of scales.}
For $0<\varepsilon<R$ define
\[
 T_{\varepsilon,R}(f,g,h)(x)=
 \int_{\varepsilon<|t|<R}f(x+t)g(x+2t)h(x+3t)\,\frac{dt}{t}.
\]
Translation invariance of $L^4$ and Holder's inequality give, for every $t$,
\[
 \|f(\cdot+t)g(\cdot+2t)h(\cdot+3t)\|_{4/3}
 \le\|f\|_4\|g\|_4\|h\|_4.
\]
Minkowski's integral inequality therefore yields
\[
 \|T_{\varepsilon,R}(f,g,h)\|_{4/3}
 \le2\log(R/\varepsilon)\|f\|_4\|g\|_4\|h\|_4.
 \tag{444.1}
\]
In particular each annulus $2^j<|t|<2^{j+1}$ has norm at most $2\log2$ times the product of the input norms. Summing those bounds reintroduces the number of annuli. Uniformity for one annulus does not give uniformity for their signed sum. Equation (444.1) is an upper estimate, so its divergence is not an example of an unbounded operator.

\paragraph{Cancellation near zero is real but costs derivatives.}
Put $F_t(x)=f(x+t)g(x+2t)h(x+3t)$. Symmetric integration replaces the integrand near zero by $(F_t-F_{-t})/t$. Differentiation in the parameter and Holder's inequality give
\[
 \|\partial_tF_t\|_{4/3}\le A,
 \qquad
 A=\|f'\|_4\|g\|_4\|h\|_4
 +2\|f\|_4\|g'\|_4\|h\|_4
 +3\|f\|_4\|g\|_4\|h'\|_4.
\]
Integrating this derivative from $-t$ to $t$ proves
\[
 \left\|\int_0^r\frac{F_t-F_{-t}}t\,dt\right\|_{4/3}
 \le2rA. \tag{444.2}
\]
Thus the principal value exists near zero for Schwartz inputs, including convergence in this norm locally in the parameter. At large $|t|$, the map $(x,t)\mapsto(x+t,x+2t)$ is invertible, so the product of the first two Schwartz factors decays rapidly in $(x,t)$; the third is bounded. This controls the tail for fixed Schwartz functions. It proves that the individual expression is meaningful, but the constants involve derivatives or decay bounds absent from the proposed inequality.

Rescaling does not remove that dependence. Under $f_a(x)=f(ax)$, and likewise for $g,h$, the operator rescales as $T(f_a,g_a,h_a)(x)=T(f,g,h)(ax)$ for $a>0$. Both sides of the desired norm inequality scale as $a^{-3/4}$. In (444.2), differentiation adds a factor $a$, compensated only by shrinking $r$. The remaining scales still require an estimate.

\paragraph{Fourier calculation and a fully bounded restricted class.}
Use $\widehat f(\xi)=\int f(x)e^{-ix\xi}\,dx$ with inverse factor $(2\pi)^{-1}$. To check the kernel rigorously, first insert $e^{-\eta|t|}$, $\eta>0$, in its symmetric principal value. For real $u$,
\[
 \operatorname{p.v.}\int_{\mathbb R}e^{-\eta|t|}e^{iut}\,\frac{dt}{t}
 =2i\int_0^\infty e^{-\eta t}\frac{\sin(ut)}t\,dt
 =2i\arctan(u/\eta).
 \tag{444.3}
\]
The last identity follows by differentiating in $u$, integrating $e^{-\eta t}\cos(ut)$, and setting the integration constant at $u=0$. As $\eta\downarrow0$ the multiplier converges to $i\pi\operatorname{sgn}u$. Consequently
\[
 T(f,g,h)(x)=\frac{i\pi}{(2\pi)^3}
 \int_{\mathbb R^3}e^{ix(\xi_1+\xi_2+\xi_3)}
 \operatorname{sgn}(\xi_1+2\xi_2+3\xi_3)
 \widehat f(\xi_1)\widehat g(\xi_2)\widehat h(\xi_3)\,d\xi.
 \tag{444.4}
\]
For Schwartz transforms, their product is integrable in the three separate frequency variables. The bounded damped multipliers permit dominated convergence. The spatial damping converges to the already justified principal value, so the two descriptions agree.

Suppose the three Fourier supports lie in compact intervals $I_1,I_2,I_3$ satisfying
\[
 \inf\{\xi_1+2\xi_2+3\xi_3:\xi_j\in I_j\}>0.
 \tag{444.5}
\]
Then the multiplier is constant on their product. Fourier inversion gives the exact identity $T(f,g,h)=i\pi fgh$, and hence
\[
 \|T(f,g,h)\|_{4/3}\le\pi\|f\|_4\|g\|_4\|h\|_4.
 \tag{444.6}
\]
The same result with a minus sign holds when the supremum of the linear form is negative. Such Schwartz inputs exist by choosing smooth compactly supported transforms. This proves a nonempty restricted class without any dependence on interval widths or distances from the singular plane, provided their product stays strictly on one side.

For general inputs, decomposing Fourier space into these regions is not an automatic proof. The cutting plane couples all three frequency variables, so its indicator does not generally factor into three independent projections with summable operator norms. A pointwise bounded multilinear multiplier also does not give the desired estimate by Plancherel. Plancherel controls a single $L^2$ transform, whereas the inputs here are in $L^4$ and three frequencies contribute to one output frequency.

\paragraph{Exact linear and quadratic modulation identities.}
For real $a,b,c$ with $a+2b+3c=0$, substitution into the original integral gives
\[
 T(e^{iax}f,e^{ibx}g,e^{icx}h)(x)
 =e^{i(a+b+c)x}T(f,g,h)(x). \tag{444.7}
\]
More strikingly, for every real $u$,
\[
 3(x+t)^2-3(x+2t)^2+(x+3t)^2=x^2,
\]
and therefore
\[
 T(e^{3iux^2}f,e^{-3iux^2}g,e^{iux^2}h)(x)
 =e^{iux^2}T(f,g,h)(x). \tag{444.8}
\]
Multiplication by these phases preserves all relevant norms and takes Schwartz functions to Schwartz functions. Thus arbitrarily large quadratic oscillation in the three factors can cancel completely in the integration parameter. An argument claiming uniform gain merely because each input oscillates rapidly would contradict (444.8). This is an obstruction to that proposed method, not an obstruction to boundedness itself.

The balanced formulation becomes especially transparent after duality. With $q\in L^4$, pair the output with $\overline q$ to obtain the four-function expression
\[
 \operatorname{p.v.}\iint
 \overline{q(x)}f(x+t)g(x+2t)h(x+3t)\,\frac{dt\,dx}{t}.
 \tag{444.9}
\]
The desired estimate is a bound for this form by the product of four $L^4$ norms. The coefficient vector $(-1,3,-3,1)$ annihilates the zeroth, first, and second powers of the four slopes $0,1,2,3$. This directly accounts for the quadratic symmetry in the form and explains why a proof must accommodate coherent phase patterns across scales.

\paragraph{Verification and outcome.}
Exact polynomial checks verify the phase cancellations and the coefficient conditions, including the sign of the output phase. The derivative and frequency-restricted estimates are proved above, not inferred from numerical quadrature of a singular kernel. The cited cancellation and curved-operator results were checked for their actual scope.

\textbf{UNRESOLVED.} We obtain the uniform bound (444.6) under the frequency-support condition and derivative-dependent control for arbitrary Schwartz data. Removing support restrictions while summing all scales with a constant independent of the inputs and truncations remains unproved; no counterexample to the balanced estimate is constructed.

\subsection{Sources}
\begin{itemize}
\item[S1]Terence Tao, Open question: boundedness of the trilinear Hilbert transform (10 May 2007).
\url{https://terrytao.wordpress.com/2007/05/10/open-question-boundedness-of-the-trilinear-hilbert-transform/}.
\textit{Formulation source.}
\item[E1]T. Tao, \emph{Cancellation for the multilinear Hilbert transform}, arXiv:1505.06479.
\url{https://arxiv.org/abs/1505.06479}.
\textit{Cancellation with truncation dependence, not the full boundedness conjecture.}
\item[E2]A. Gaitan and V. Lie, \emph{On the curved Trilinear Hilbert transform}, arXiv:2308.10706.
\url{https://arxiv.org/abs/2308.10706}.
\textit{Curved shifts are outside the fixed linear-slope target; consulted 27 September 2026.}
\end{itemize}
\end{document}
